\chapter{Convex bodies}\label{ch:convex}

\begin{center}\itshape
They bought as much land as they could enclose with the hide of a bull.
\upshape(paraphrase of Virgil, \emph{Aeneid}, Book~I, lines 367--368)
\end{center}

Queen Dido, fleeing Tyre, landed on the coast of Africa and was allowed to buy as much land as a bull's hide could enclose. She cut the hide into thin strips,
tied them into a long cord, and enclosed a whole hill, on which Carthage was built. Which shape encloses the most area for a cord of given length? Every
schoolchild ``knows'' that it is the circle, and the ancient Greeks knew it too, but a complete proof had to wait for the nineteenth century. Tome~4
(Chapter~6) gives Hurwitz's proof with Fourier series. Here we prove it in every dimension, as a corollary of a beautiful inequality about adding
shapes, the Brunn--Minkowski inequality.

Convex sets are the natural home of such questions. We meet three classics on the way: Helly's theorem, a combinatorial jewel; Buffon's needle, the
first problem of geometric probability, which measures $\pi$ by throwing sticks on a floor; and the curves of constant width, which roll between two parallel
lines like wheels without being circles.

\begin{think}
A needle of length $\ell<1$ is dropped at random on a floor ruled with parallel lines at distance $1$. Which do you expect to be more likely: that it
crosses a line, or that it does not, when $\ell=\frac34$? Does the answer depend on the shape of the needle if you bend it?
\end{think}

\section{Convex sets and Helly's theorem}

A set $C\subset\mathbb R^n$ is \emph{convex} if it contains the segment between any two of its points. Intersections of convex sets are convex. The
\emph{convex hull} $\operatorname{conv}(S)$ of a set $S$ is the smallest convex set containing it; it consists of the convex combinations
$\sum t_ix_i$ ($x_i\in S$, $t_i\ge0$, $\sum t_i=1$) of finitely many points of $S$.

\begin{lemma}[Radon, 1921]\label{lem:radon}
Any $m\ge n+2$ points $x_1,\dots,x_m$ of $\mathbb R^n$ can be split into two sets $I$, $J$ whose convex hulls meet.
\end{lemma}

\begin{proof}
The $m$ vectors $(x_i,1)\in\mathbb R^{n+1}$ are linearly dependent since $m>n+1$: there are $\lambda_i$, not all zero, with $\sum\lambda_ix_i=0$ and $\sum\lambda_i=0$.
Let $I=\{i:\lambda_i>0\}$ and $J$ the rest; both contain an index with $\lambda_i\ne0$, since the $\lambda_i$ add up to $0$. With
$s=\sum_{i\in I}\lambda_i=\sum_{j\in J}(-\lambda_j)>0$, the point $p=\sum_{i\in I}\frac{\lambda_i}sx_i=\sum_{j\in J}\frac{-\lambda_j}sx_j$ lies in both convex hulls.
\end{proof}

\begin{theorem}[Helly]\label{thm:helly}
Let $C_1,\dots,C_m$ be convex sets in $\mathbb R^n$, $m\ge n+1$. If every $n+1$ of them have a common point, then all of them have a common point.
\end{theorem}

\begin{proof}
By induction on $m$; for $m=n+1$ there is nothing to prove. Let $m\ge n+2$. By induction, any $m-1$ of the sets have a common point: choose
$x_i\in\bigcap_{j\ne i}C_j$ for each $i$. By Radon's lemma, split the indices into $I$ and $J$ with a point $p\in\operatorname{conv}\{x_i:i\in I\}\cap
\operatorname{conv}\{x_j:j\in J\}$. Take any $k$. If $k\notin I$, then every $x_i$ with $i\in I$ lies in $C_k$ (because $i\ne k$), so their convex hull, and $p$, lie in
the convex set $C_k$. If $k\in I$, then $k\notin J$ and the same argument with $J$ applies. So $p$ is in every $C_k$.
\end{proof}

Eduard Helly found the theorem in 1913 but published it only in 1923; in between he spent years as a prisoner of war in Siberia, and Johann Radon
published the first proof in 1921 (see Danzer, Gr\"unbaum and Klee in the sources). The number $n+1$ is sharp (Exercise~\ref{ex:hellysharp}), and the
convexity is essential: two disjoint convex pieces do not make a convex set.

\section{Buffon's needle}

In 1733 Georges-Louis Leclerc, comte de Buffon, presented to the Paris Academy a problem about a game: a coin or a stick is thrown on a tiled floor; what is the
chance that it falls across a joint? He published the solution in 1777. It was the birth of \emph{geometric probability}. We model the floor by the lines
$y\in\mathbb Z$ and throw a rigid figure: we rotate it by an angle $\theta$ uniform in $[0,2\pi)$ and shift it vertically by $u$ uniform in $[0,1)$, independent
(horizontal shifts do not matter). Let $N$ be the number of points where the figure meets the lines.

\begin{lemma}\label{lem:randominterval}
For $h\ge0$ and $u$ uniform in $[0,1)$, the interval $[u,u+h)$ contains on average exactly $h$ integers.
\end{lemma}

\begin{proof}
Write $h=k+r$ with $k$ an integer and $0\le r<1$. The interval $[u,u+k)$ contains exactly $k$ integers, and $[u+k,u+k+r)$ contains one integer if
$u+k+r>\lceil u+k\rceil$, that is, with probability $r$ (for $u\in(1-r,1)$, or $u=0$), and none otherwise. The mean is $k+r=h$.
\end{proof}

\begin{theorem}[Buffon, Barbier]\label{thm:buffon}
For a polygonal curve of total length $L$, thrown as above, $\mathbb E[N]=\frac{2L}\pi$. In particular a needle of length $\ell\le1$ crosses a line with
probability $\frac{2\ell}\pi$.
\end{theorem}

\begin{proof}
$N$ is the sum of the numbers of crossings of the segments of the curve (up to events of probability zero, where a vertex lies on a line), so by
linearity of expectation it suffices to treat one segment, of length $\ell$. Rotated by $\theta$, it occupies a vertical extent $h=\ell|\sin\theta|$, and for
almost every $u$ it meets the lines at the integers of a half-open interval of length $h$, shifted by $u$. By Lemma~\ref{lem:randominterval}, the expected
number of crossings given $\theta$ is $\ell|\sin\theta|$, so
\[
\mathbb E[N]=\frac1{2\pi}\int_0^{2\pi}\ell|\sin\theta|\,d\theta=\frac{4\ell}{2\pi}=\frac{2\ell}\pi.
\]
A needle of length $\ell\le1$ meets at most one line (almost surely), so $N\in\{0,1\}$ and $\Pr(N=1)=\mathbb E[N]$.
\end{proof}

The striking part, due to Joseph-\'Emile Barbier (1860), is that the shape does not matter: a bent needle, a ``noodle'', crosses on average as often as a
straight one of the same length. A circle of diameter $1$, of length $\pi$, always crosses exactly twice; that alone determines the constant
$\frac2\pi$, without any integral.

\begin{figure}[ht]
\centering
\fig{ch11-buffon}
\caption{Forty needles of length $0.8$ thrown on lines at distance $1$. Here $20$ of them cross a line (in red); the theorem predicts a proportion
$\frac{1.6}\pi\approx0.51$.}
\label{fig:buffon}
\end{figure}

So $\pi\approx\frac{2\ell\cdot(\text{number of throws})}{\text{number of crossings}}$. In 1901 the Italian mathematician Mario Lazzarini reported $3408$ throws of a
needle with $\ell=\frac56$, and the estimate $\frac{355}{113}$, correct to six decimals. Too good to be true: Lee Badger showed in 1994 how the numbers were
almost certainly chosen to produce this fraction (Exercise~\ref{ex:lazzarini}). An honest experiment gains one correct digit for every hundredfold increase in
the number of throws, by the central limit theorem of Tome~3.

\section{Width and curves of constant width}

Let $K$ be a compact convex set in the plane. Its \emph{width} in the direction $\theta$, $w_K(\theta)$, is the length of its orthogonal projection onto the
line of direction $\theta$: the distance between the two support lines of $K$ perpendicular to that line. Its \emph{perimeter} $L(K)$ is the length of its
boundary, the supremum of the perimeters of the convex polygons inscribed in it.

\begin{theorem}[Cauchy's formula]\label{thm:cauchy}
For a compact convex set $K$ in the plane, $L(K)=\int_0^\pi w_K(\theta)\,d\theta$. So the perimeter is $\pi$ times the mean width.
\end{theorem}

\begin{proof}
First let $K$ be a convex polygon with sides of lengths $\ell_i$ and directions $\theta_i$. Projecting onto the line of direction $\theta$, the boundary covers the
projection of $K$ twice (once going, once coming back), so $w_K(\theta)=\frac12\sum_i\ell_i|\cos(\theta-\theta_i)|$. Since
$\int_0^\pi|\cos(\theta-\theta_i)|\,d\theta=2$, we get $\int_0^\pi w_K=\sum_i\ell_i=L(K)$.

In general, choose inscribed convex polygons $K_j$ whose vertices become dense on the boundary of $K$, each containing the previous one, with
$L(K_j)\to L(K)$ (adding a vertex can only increase the perimeter, by the triangle inequality). Every point of $K$ is within $\varepsilon_j$ of $K_j$, with
$\varepsilon_j\to0$, so $w_{K_j}\le w_K\le w_{K_j}+2\varepsilon_j$. Passing to the limit in $L(K_j)=\int_0^\pi w_{K_j}$ gives the formula.
\end{proof}

A convex set has \emph{constant width} $w$ if $w_K(\theta)=w$ for all $\theta$. The disk of diameter $w$ is one. The \emph{Reuleaux triangle} is another:
take an equilateral triangle of side $w$ and replace each side by the arc of the circle of radius $w$ centred at the opposite vertex (Figure~\ref{fig:reuleaux}).
Between two parallel support lines, one touches a vertex and the other the opposite arc, at distance $w$.

\begin{corollary}[Barbier, 1860]\label{cor:barbier}
Every convex set of constant width $w$ has perimeter $\pi w$.
\end{corollary}

\begin{proof}
By Cauchy's formula, $L=\int_0^\pi w\,d\theta=\pi w$.
\end{proof}

\begin{figure}[ht]
\centering
\fig{ch11-reuleaux}
\caption{A Reuleaux triangle and a disk of constant width $1$, each turning inside a unit square while touching its four sides. Both have perimeter $\pi$.}
\label{fig:reuleaux}
\end{figure}

Curves of constant width are named after Franz Reuleaux, an engineer who studied mechanisms in the 1870s. They make honest wheels for a rolling platform on
rollers (the platform stays level), and the Reuleaux triangle is the basis of drills that cut almost square holes. Among all sets of constant width $w$, the
Reuleaux triangle has the smallest area (the Blaschke--Lebesgue theorem, a remark here) and the disk the largest (Exercise~\ref{ex:constwidthdisk}).

\section{The Brunn--Minkowski inequality}

The \emph{Minkowski sum} of two sets is $A+B=\{a+b:a\in A,\ b\in B\}$. For instance, $K+\varepsilon\mathbb B$, with $\mathbb B$ the unit ball, is the set of
points at distance at most $\varepsilon$ from $K$.

\begin{figure}[ht]
\centering
\fig{ch11-minkowski}
\caption{The Minkowski sum of a triangle $A$ (area $1.4$) and a square $B$ (area $0.49$) is a hexagon of area $4.27$, and
$\sqrt{4.27}\approx2.07\ge\sqrt{1.4}+\sqrt{0.49}\approx1.88$.}
\label{fig:minkowskisum}
\end{figure}

\begin{theorem}[Brunn--Minkowski]\label{thm:bm}
For nonempty compact sets $A,B\subset\mathbb R^n$,
\[
\operatorname{vol}(A+B)^{1/n}\ge\operatorname{vol}(A)^{1/n}+\operatorname{vol}(B)^{1/n}.
\]
\end{theorem}

\begin{proof}
A \emph{box} is a product of closed intervals of positive length. We first prove the inequality when $A$ and $B$ are finite unions of boxes with disjoint
interiors, by induction on the total number $k$ of boxes. Write $|X|$ for $\operatorname{vol}(X)$.

\emph{Two boxes.} If $A$ and $B$ are boxes with sides $a_i$ and $b_i$, then $A+B$ is a box with sides $a_i+b_i$, and by the inequality between arithmetic and geometric
means,
\[
\prod_i\Bigl(\frac{a_i}{a_i+b_i}\Bigr)^{1/n}+\prod_i\Bigl(\frac{b_i}{a_i+b_i}\Bigr)^{1/n}\le\frac1n\sum_i\frac{a_i}{a_i+b_i}+\frac1n\sum_i\frac{b_i}{a_i+b_i}=1,
\]
which is $|A|^{1/n}+|B|^{1/n}\le|A+B|^{1/n}$.

\emph{Induction.} Let $k\ge3$, and say $A$ has at least two boxes. Two boxes with disjoint interiors are separated by a hyperplane $x_j=t$ (in some coordinate their
intervals overlap at most at an endpoint). Let $A^\pm=A\cap\{\pm(x_j-t)\ge0\}$: cutting each box of $A$ by the hyperplane and discarding pieces of volume zero, each is a union of boxes with disjoint
interiors, and each has fewer boxes than $A$, because one of the two separated boxes lies entirely on the other side. Let $\lambda=|A^+|/|A|\in(0,1)$. The function $s\mapsto|B\cap\{x_j\ge s\}|$ is
continuous and decreases from $|B|$ to $0$, so there is $s$ with $|B^+|=\lambda|B|$ and $|B^-|=(1-\lambda)|B|$, where $B^\pm=B\cap\{\pm(x_j-s)\ge0\}$; these have
at most as many boxes as $B$. Now $A^++B^+$ lies in $\{x_j\ge t+s\}$ and $A^-+B^-$ in $\{x_j\le t+s\}$, so they overlap only in a set of volume zero, and both lie in
$A+B$. By the induction hypothesis,
\[
|A+B|\ge|A^++B^+|+|A^-+B^-|\ge\bigl(|A^+|^{\frac1n}+|B^+|^{\frac1n}\bigr)^n+\bigl(|A^-|^{\frac1n}+|B^-|^{\frac1n}\bigr)^n.
\]
Since $|A^+|=\lambda|A|$ and $|B^+|=\lambda|B|$, the first term is $\lambda\bigl(|A|^{\frac1n}+|B|^{\frac1n}\bigr)^n$, the second is $(1-\lambda)$ times the same, and the
sum is $\bigl(|A|^{\frac1n}+|B|^{\frac1n}\bigr)^n$.

\emph{Compact sets.} For $j\ge1$, let $A_j$ be the union of the closed cubes of the grid $2^{-j}\mathbb Z^n$ that meet $A$, and similarly $B_j$. These are finite
unions of boxes with disjoint interiors, they decrease with $j$, and their intersection is $A$ (resp.\ $B$), since $A$ is closed; so
$|A_j|\to|A|$ and $|B_j|\to|B|$. Every point of $A_j$ is within $\delta_j=\sqrt n\,2^{-j}$ of $A$, so $A_j+B_j\subset(A+B)+2\delta_j\mathbb B$, a decreasing family of
sets whose intersection is the compact set $A+B$. Hence
\[
|A+B|=\lim_j|(A+B)+2\delta_j\mathbb B|\ge\limsup_j|A_j+B_j|\ge\lim_j\bigl(|A_j|^{\frac1n}+|B_j|^{\frac1n}\bigr)^n=\bigl(|A|^{\frac1n}+|B|^{\frac1n}\bigr)^n.
\]
\end{proof}

Hermann Brunn proved the inequality for convex bodies in 1887 and Minkowski completed it; the slicing proof above is due to Hadwiger and Ohmann (1956).
Richard Gardner's survey in the sources tells how this one inequality sits at the crossroads of geometry, analysis and probability.

\begin{corollary}[isoperimetric inequality]\label{cor:isoperimetric}
For a compact set $K\subset\mathbb R^n$ with $|K|>0$, define its surface area as $S(K)=\liminf_{\varepsilon\to0^+}\frac{|K+\varepsilon\mathbb B|-|K|}\varepsilon$. Then
\[
S(K)\ge n\,|\mathbb B|^{1/n}|K|^{(n-1)/n},
\]
with equality for balls. In the plane: $L^2\ge4\pi A$.
\end{corollary}

\begin{proof}
By Brunn--Minkowski, $|K+\varepsilon\mathbb B|\ge\bigl(|K|^{1/n}+\varepsilon|\mathbb B|^{1/n}\bigr)^n=|K|+n\varepsilon|K|^{(n-1)/n}|\mathbb B|^{1/n}+O(\varepsilon^2)$.
Subtract $|K|$, divide by $\varepsilon$ and let $\varepsilon\to0$. For a ball $K=r\mathbb B$, $K+\varepsilon\mathbb B=(r+\varepsilon)\mathbb B$ and both sides equal
$nr^{n-1}|\mathbb B|$. In the plane, $|\mathbb B|=\pi$, so $S\ge2\sqrt{\pi A}$, that is $S^2\ge4\pi A$; for a convex set $S$ is the perimeter $L$
(Exercise~\ref{ex:steiner}).
\end{proof}

So among all bodies of a given volume, the ball has the least surface area: soap bubbles are round. Steiner gave five ingenious arguments in 1838 that if a
best shape exists, it is the disk; Weierstrass pointed out that existence is the whole difficulty, and the first complete proofs came later in the
nineteenth century. In the proof above existence never enters: the inequality holds directly for every compact set.

\section{Exercises}

\begin{exercise}[$\star$]\label{ex:helly1d}
Prove Helly's theorem in dimension $1$ directly: if finitely many intervals meet pairwise, they all have a common point.
\end{exercise}

\begin{exercise}[$\star$]\label{ex:hellysharp}
Find $n+1$ convex sets in $\mathbb R^n$ such that any $n$ of them meet but the $n+1$ do not.
\end{exercise}

\begin{exercise}[$\star\star$]\label{ex:hellydisk}
Let $S$ be a finite set of points in the plane such that any three of them lie in some disk of radius $1$. Show that all of $S$ lies in a disk of radius $1$.
\end{exercise}

\begin{exercise}[$\star\star$]\label{ex:caratheodory}
(Carath\'eodory) Show that every point of $\operatorname{conv}(S)$, for $S\subset\mathbb R^n$, is a convex combination of at most $n+1$ points of $S$.
\end{exercise}

\begin{exercise}[$\star$]\label{ex:lazzarini}
Lazzarini threw a needle of length $\ell=\frac56$ (lines at distance $1$) $3408$ times. How many crossings give the estimate $\frac{355}{113}$? Why is the number
$3408$ suspicious?
\end{exercise}

\begin{exercise}[$\star$]\label{ex:noodle}
Compute $\mathbb E[N]$ for a square of side $a$ and for a circle of radius $r$, thrown on lines at distance $1$. For which $a$ is a square of side $a$ certain to meet
at least one line?
\end{exercise}

\begin{exercise}[$\star$]\label{ex:reuleauxarea}
Compute the area of a Reuleaux triangle of width $1$, and compare with the disk of diameter $1$.
\end{exercise}

\begin{exercise}[$\star\star$]\label{ex:minkowskiwidth}
Show that $w_{K+L}=w_K+w_L$ for convex $K$, $L$ in the plane, and deduce that $L(K+L)=L(K)+L(L)$. What is the perimeter of $K+\varepsilon\mathbb B$?
\end{exercise}

\begin{exercise}[$\star\star$]\label{ex:steiner}
(Steiner's formula) For a convex polygon $K$ of area $A$ and perimeter $L$, show that $|K+\varepsilon\mathbb B|=A+L\varepsilon+\pi\varepsilon^2$. Deduce that the surface
area of Corollary~\ref{cor:isoperimetric} is the perimeter.
\end{exercise}

\begin{exercise}[$\star\star$]\label{ex:constwidthdisk}
Using the isoperimetric inequality and Barbier's theorem, show that among convex sets of constant width $w$, the disk has the largest area.
\end{exercise}

\begin{exercise}[$\star$]\label{ex:bmone}
Prove the Brunn--Minkowski inequality in dimension $1$ directly: $|A+B|\ge|A|+|B|$ for nonempty compact $A,B\subset\mathbb R$.
\end{exercise}

\begin{exercise}[$\star\star$]\label{ex:bmconcave}
Show that for convex $A$ and $0\le\lambda\le1$, $(1-\lambda)A+\lambda A=A$, and deduce from Theorem~\ref{thm:bm} that for compact $A,B$ the function
$\lambda\mapsto|(1-\lambda)A+\lambda B|^{1/n}$ is at least $(1-\lambda)|A|^{1/n}+\lambda|B|^{1/n}$. Apply it to the slices of a convex body: if $K\subset\mathbb R^3$ is convex,
the square root of the area of the slice $K\cap\{z=t\}$ is a concave function of $t$ where it is positive.
\end{exercise}

\section{Hints}

\begin{hint}{ex:helly1d}
Compare the largest left endpoint with the smallest right endpoint.
\end{hint}
\begin{hint}{ex:hellysharp}
Use the $n+1$ facets of a simplex, or the half-spaces bounded by them.
\end{hint}
\begin{hint}{ex:hellydisk}
Apply Helly to the disks of radius $1$ centred at the points of $S$.
\end{hint}
\begin{hint}{ex:caratheodory}
If $x=\sum_{i=1}^mt_ix_i$ with $m\ge n+2$, use an affine dependence $\sum\lambda_ix_i=0$, $\sum\lambda_i=0$ to remove one point.
\end{hint}
\begin{hint}{ex:lazzarini}
The estimate is $\frac{2\ell T}C$ for $T$ throws and $C$ crossings.
\end{hint}
\begin{hint}{ex:noodle}
Use Theorem~\ref{thm:buffon} for the square; for the circle, use the circle of diameter $1$.
\end{hint}
\begin{hint}{ex:reuleauxarea}
The triangle plus three circular segments.
\end{hint}
\begin{hint}{ex:minkowskiwidth}
The projection of $K+L$ is the sum of the projections.
\end{hint}
\begin{hint}{ex:steiner}
Split $K+\varepsilon\mathbb B$ into $K$, rectangles on the sides, and sectors at the vertices.
\end{hint}
\begin{hint}{ex:constwidthdisk}
$4\pi A\le L^2=\pi^2w^2$.
\end{hint}
\begin{hint}{ex:bmone}
Translate so that $\max A=0=\min B$.
\end{hint}
\begin{hint}{ex:bmconcave}
Apply Theorem~\ref{thm:bm} to $(1-\lambda)A$ and $\lambda B$; for the slices, check that $(1-\lambda)K_s+\lambda K_t\subset K_{(1-\lambda)s+\lambda t}$.
\end{hint}

\section{Solutions}

\begin{solution}{ex:helly1d}
\idea Intervals meeting pairwise have every left end to the left of every right end.
\writeup Let the intervals be $[a_i,b_i]$ (the open or half-open cases are similar). For all $i,j$, the intervals $i$ and $j$ meet, so $a_i\le b_j$. Hence
$a=\max_ia_i\le\min_jb_j=b$, and every point of $[a,b]$, for instance $a$, lies in all of them.
\end{solution}

\begin{solution}{ex:hellysharp}
\idea A simplex is the intersection of $n+1$ half-spaces, and any $n$ of its facet hyperplanes meet at a vertex.
\writeup Let $v_0,\dots,v_n$ be the vertices of a simplex and $F_i$ the facet opposite $v_i$, a convex set. Any $n$ facets, all but $F_i$, contain the vertex $v_i$.
But no point lies on all $n+1$ facets: a point of the simplex has barycentric coordinates $t_0,\dots,t_n\ge0$ with sum $1$, it lies on $F_i$ exactly when $t_i=0$,
and they cannot all vanish.
\end{solution}

\begin{solution}{ex:hellydisk}
\idea ``Some disk of radius $1$ contains these points'' means ``the disks of radius $1$ around these points have a common point''.
\writeup A disk of radius $1$ centred at $c$ contains $p$ if and only if $c$ lies in the disk $D_p$ of radius $1$ centred at $p$. The hypothesis says that any three of the
convex sets $D_p$, $p\in S$, have a common point. By Helly's theorem in the plane, all of them have a common point $c$, and the disk of radius $1$ centred at $c$
contains $S$. (If $S$ has fewer than three points, the hypothesis applies to $S$ itself.)
\end{solution}

\begin{solution}{ex:caratheodory}
\idea Radon's trick, used to shift weight off one point.
\writeup Let $x=\sum_{i=1}^mt_ix_i$ be a convex combination with $m$ minimal, so all $t_i>0$. If $m\ge n+2$, take $\lambda_i$, not all zero, with $\sum\lambda_ix_i=0$ and
$\sum\lambda_i=0$ (as in Lemma~\ref{lem:radon}); some $\lambda_i>0$. Let $c=\min\{t_i/\lambda_i:\lambda_i>0\}$. Then $x=\sum(t_i-c\lambda_i)x_i$, the coefficients are
$\ge0$, they add up to $1$, and at least one vanishes: a shorter combination, contradicting minimality. So $m\le n+1$.
\end{solution}

\begin{solution}{ex:lazzarini}
\idea Work backwards from the fraction.
\writeup The estimate is $\frac{2\cdot\frac56\cdot3408}C=\frac{5680}C$. It equals $\frac{355}{113}$ when $C=\frac{5680\cdot113}{355}=1808$. Now $\frac{2\ell}\pi=\frac5{3\pi}$, and
$\frac53\cdot\frac{113}{355}=\frac{113}{213}$: to get $\frac{355}{113}$ one needs a number of throws that is a multiple of $213$ and exactly $\frac{113}{213}$ of them
crossing. Indeed $3408=16\cdot213$ and $1808=16\cdot113$. An experimenter who stops exactly when the ratio hits a famous fraction, or who reports a convenient
number of throws, can produce any accuracy; Badger's article analyses how unlikely an honest $\frac{355}{113}$ would be.
\end{solution}

\begin{solution}{ex:noodle}
\idea Only the length matters.
\writeup The square has perimeter $4a$, so $\mathbb E[N]=\frac{8a}\pi$. The circle has length $2\pi r$, so $\mathbb E[N]=4r$ (for $r=\frac12$ it crosses exactly twice, as
it should). A square of side $a$ always meets a line when its vertical extent, between $a$ and $a\sqrt2$ depending on the angle, is at least $1$ for every angle,
that is, when $a\ge1$.
\end{solution}

\begin{solution}{ex:reuleauxarea}
\idea Triangle plus three caps.
\writeup The equilateral triangle of side $1$ has area $\frac{\sqrt3}4$. Each arc bounds, together with the opposite side, a circular segment: a sector of angle
$\frac\pi3$ and radius $1$, of area $\frac\pi6$, minus the triangle, of area $\frac{\sqrt3}4$. So the area is
$\frac{\sqrt3}4+3\bigl(\frac\pi6-\frac{\sqrt3}4\bigr)=\frac{\pi-\sqrt3}2\approx0.705$, against $\frac\pi4\approx0.785$ for the disk of diameter $1$.
\end{solution}

\begin{solution}{ex:minkowskiwidth}
\idea Projections commute with Minkowski sums.
\writeup Let $\pi_\theta$ be the orthogonal projection onto the line of direction $\theta$. It is linear, so $\pi_\theta(K+L)=\pi_\theta(K)+\pi_\theta(L)$, and the sum of two
segments of lengths $w_K(\theta)$ and $w_L(\theta)$ is a segment of length $w_K(\theta)+w_L(\theta)$. Integrating over $\theta$ with Cauchy's formula,
$L(K+L)=L(K)+L(L)$ ($K+L$ is convex). In particular $L(K+\varepsilon\mathbb B)=L(K)+2\pi\varepsilon$.
\end{solution}

\begin{solution}{ex:steiner}
\idea Points at distance at most $\varepsilon$ from a convex polygon are in the polygon, beside a side, or near a vertex.
\writeup Let $K$ have sides $\ell_1,\dots,\ell_m$ and exterior angles $\phi_1,\dots,\phi_m$, which add up to $2\pi$. The set $K+\varepsilon\mathbb B$ is the union, with
disjoint interiors, of $K$, of a rectangle $\ell_i\times\varepsilon$ on each side, and of a circular sector of radius $\varepsilon$ and angle $\phi_i$ at each vertex (the
nearest point of $K$ to a point outside is on a side, or is a vertex). Hence $|K+\varepsilon\mathbb B|=A+\varepsilon\sum\ell_i+\frac{\varepsilon^2}2\sum\phi_i=A+L\varepsilon+\pi\varepsilon^2$,
and $\frac{|K+\varepsilon\mathbb B|-A}\varepsilon\to L$. For a general convex set, approximating by polygons as in Cauchy's formula gives the same result.
\end{solution}

\begin{solution}{ex:constwidthdisk}
\idea Fix the perimeter, then maximise the area.
\writeup By Barbier's theorem, every convex set of constant width $w$ has perimeter $\pi w$. By the isoperimetric inequality, its area is at most
$\frac{L^2}{4\pi}=\frac{\pi w^2}4$, the area of the disk of diameter $w$.
\end{solution}

\begin{solution}{ex:bmone}
\idea Place $A$ to the left of $B$, touching.
\writeup Translations do not change the measures, so assume $\max A=0=\min B$. Then $A=A+0\subset A+B$ and $B=0+B\subset A+B$, and $A\cap B=\{0\}$. So
$|A+B|\ge|A\cup B|=|A|+|B|$.
\end{solution}

\begin{solution}{ex:bmconcave}
\idea Convex combinations of a convex set stay in it; then Brunn--Minkowski with scaled sets.
\writeup For convex $A$, $(1-\lambda)a+\lambda a'\in A$, so $(1-\lambda)A+\lambda A\subset A$, and $a=(1-\lambda)a+\lambda a$ gives the reverse inclusion. For compact $A,B$,
Theorem~\ref{thm:bm} applied to $(1-\lambda)A$ and $\lambda B$, whose volumes are $(1-\lambda)^n|A|$ and $\lambda^n|B|$, gives
$|(1-\lambda)A+\lambda B|^{1/n}\ge(1-\lambda)|A|^{1/n}+\lambda|B|^{1/n}$. For the slices $K_t=\{(x,y):(x,y,t)\in K\}$: if $p\in K_s$ and $q\in K_t$, then
$(1-\lambda)(p,s)+\lambda(q,t)\in K$ by convexity, so $(1-\lambda)K_s+\lambda K_t\subset K_{(1-\lambda)s+\lambda t}$. Brunn--Minkowski in the plane ($n=2$) then gives
$\sqrt{|K_{(1-\lambda)s+\lambda t}|}\ge(1-\lambda)\sqrt{|K_s|}+\lambda\sqrt{|K_t|}$: concavity. This is Brunn's original form of the inequality.
\end{solution}

\section*{Sources and further reading}

\begin{itemize}
\item Virgil, \emph{Aeneid}, Book~I, 365--368 (Dido and the bull's hide).
\item G.-L.~Leclerc de Buffon, \emph{Essai d'arithm\'etique morale}, in \emph{Histoire naturelle}, Suppl\'ement~4, 1777; his 1733 memoir in the \emph{Histoire
de l'Acad\'emie Royale des Sciences}; J.-\'E.~Barbier, \emph{Note sur le probl\`eme de l'aiguille et le jeu du joint couvert}, J.\ Math.\ Pures Appl.\ 5 (1860).
\item L.~Badger, \emph{Lazzarini's lucky approximation of $\pi$}, Mathematics Magazine 67 (1994), 83--91.
\item L.~Danzer, B.~Gr\"unbaum and V.~Klee, \emph{Helly's theorem and its relatives}, in \emph{Convexity}, Proc.\ Symp.\ Pure Math.\ 7, AMS, 1963 (history of
Helly's and Radon's theorems).
\item R.~J.~Gardner, \emph{The Brunn--Minkowski inequality}, Bulletin of the AMS 39 (2002), 355--405.
\item V.~Bl\aa sj\"o, \emph{The isoperimetric problem}, American Mathematical Monthly 112 (2005), 526--566 (Dido, Steiner, Weierstrass and Hurwitz).
\item D.~A.~Klain and G.-C.~Rota, \emph{Introduction to Geometric Probability}, Cambridge University Press, 1997 (Buffon, Cauchy and Crofton formulas).
\end{itemize}
